
% ============================================================================
% DOCUMENT 10: Physical AI Oncology Clinical Trial Site Activation and Standard Operating Procedures
% ============================================================================

\part{Physical AI Oncology Clinical Trial Site Activation and Standard Operating Procedures}
\setcounter{section}{0}

\section{Site Activation Checklist}

\subsection{Pre-Activation Requirements}

The following requirements shall be completed before the first patient is
enrolled at a Physical AI oncology clinical trial site:

\begin{enumerate}[label=(\alph*)]
\item \textbf{Regulatory Approvals:}
  \begin{enumerate}[label=(\arabic*)]
  \item State authorization certificate issued by CDPH (per SB 1042).
  \item San Francisco conditional use permit and health care facility operating
    license (or equivalent local permits for non-San Francisco sites).
  \item IRB approval for the Physical AI oncology trial protocol.
  \item IND clearance with Physical AI System Description accepted by
    FDA~\cite{cfr312adapt}.
  \item All Physical AI systems registered with CDPH with current USL scores
    meeting minimum thresholds~\cite{usl2026}.
  \end{enumerate}

\item \textbf{Facility Readiness:}
  \begin{enumerate}[label=(\arabic*)]
  \item Building occupancy permit issued.
  \item All Physical AI systems installed, commissioned, and verified.
  \item Phase 0 simulation validation completed for all system types with
    cross-framework consistency documented~\cite{iche6r3adapt}.
  \item Network infrastructure operational with MCP-PAI five-server topology
    configured to the required conformance level.
  \item Power systems tested including UPS and generator transfer.
  \item HVAC systems verified for all clinical zones.
  \item Security systems operational with access control zones configured.
  \end{enumerate}

\item \textbf{Staffing Readiness:}
  \begin{enumerate}[label=(\arabic*)]
  \item All site safety officers trained and CDPH-certified.
  \item On-call emergency physician and pharmacist agreements executed.
  \item Maintenance technician assigned and trained.
  \item Cybersecurity operations specialist on-call roster established.
  \item All staff completed competency assessments for Physical AI systems.
  \end{enumerate}

\item \textbf{Documentation:}
  \begin{enumerate}[label=(\arabic*)]
  \item Standard operating procedures finalized for all 10 robot types.
  \item Patient-facing robot instructions available in English and
    the most common languages spoken in the service
    area~\cite{patientinstructions2026}.
  \item Informed consent documents approved by IRB with Physical AI
    disclosures~\cite{cfr50adapt}.
  \item Emergency response plan posted and distributed.
  \item Adverse event reporting workflows tested.
  \end{enumerate}
\end{enumerate}

\section{Shift Management}

\subsection{Three-Shift Structure}

\begin{enumerate}[label=(\alph*)]
\item The site shall operate on a three-shift schedule covering 24 continuous
  hours:
  \begin{enumerate}[label=(\arabic*)]
  \item Day shift: 06:00 to 14:00.
  \item Evening shift: 14:00 to 22:00.
  \item Night shift: 22:00 to 06:00.
  \end{enumerate}

\item Each shift shall be staffed by a minimum of two site safety officers, with
  a third officer added during peak demand periods (typically 08:00 to 16:00
  based on the simulation's patient arrival patterns showing 15 arrivals
  during hour 09)~\cite{newtrial2026}.

\item Shift handoff procedures shall include a structured briefing covering:
  \begin{enumerate}[label=(\arabic*)]
  \item Current patient census and active procedures.
  \item Physical AI system status (active, standby, maintenance) for all 29
    instances.
  \item Any open adverse events or safety concerns.
  \item Scheduled maintenance windows (such as the SURG-01 22:15 to 01:59
    maintenance demonstrated in the simulation)~\cite{newtrial2026}.
  \item Pending patient arrivals and departures.
  \end{enumerate}
\end{enumerate}

\subsection{Shift Performance Metrics}

Each shift shall track and report the following metrics, consistent with the
hourly data structure validated in the 24-hour simulation:

\begin{enumerate}[label=(\alph*)]
\item Patient arrivals, procedures completed, and departures during the shift.
\item Physical AI system utilization rates per type (the simulation averaged
  8.1 of 29 active, with peak of 21 of 29 active)~\cite{newtrial2026}.
\item Average patient wait time (target: 8 minutes or less).
\item Adverse events with grade, resolution time, and root cause.
\item PSL score changes per robot type per regulatory dimension (maximum
  0.3 change per hour per dimension)~\cite{pslframework2026}.
\item Any emergency stop activations or manual overrides.
\end{enumerate}

\section{Patient Processing Workflow}

\subsection{Arrival to Procedure}

\begin{enumerate}[label=(\alph*)]
\item Patient arrival and check-in via automated kiosk (estimated 2 minutes).
\item Identity verification and informed consent confirmation. Physical AI
  disclosures reviewed for any system changes since enrollment.
\item Pre-procedure safety matrix completed, including verification that all
  assigned Physical AI systems meet USL thresholds and have passed daily
  calibration checks~\cite{cfr50adapt}.
\item Patient escorted or directed to the appropriate clinical wing based on
  scheduled procedure type.
\item Patient-facing robot instructions reviewed with the patient for the
  specific systems involved in their procedure. For example, surgical robot
  patients receive the three-step surgical robot instruction sheet covering
  anesthesia induction, robot-assisted procedure, and post-operative
  recovery~\cite{patientinstructions2026}.
\item Procedure initiated under Physical AI system autonomous control with
  site safety officer oversight.
\end{enumerate}

\subsection{Procedure to Departure}

\begin{enumerate}[label=(\alph*)]
\item Procedure completed; Physical AI system performs automated post-procedure
  protocols (retracting instruments, lowering couches, generating procedure
  reports).
\item Patient transferred to recovery area; recovery duration determined by
  procedure type (2 minutes for imaging to 60 minutes for surgical or
  steerable needle procedures)~\cite{patientinstructions2026}.
\item Post-procedure vital sign monitoring by Physical AI systems until
  discharge criteria met.
\item Discharge documentation generated automatically; patient receives
  procedure summary and follow-up instructions.
\item Patient exits via reception area; transportation coordination if needed.
\end{enumerate}

\section{Physical AI System Daily Operations}

\subsection{Daily System Checks}

\begin{enumerate}[label=(\alph*)]
\item Each Physical AI system shall undergo an automated daily verification
  check before the first patient procedure, including:
  \begin{enumerate}[label=(\arabic*)]
  \item Mechanical calibration verification (positioning accuracy, force
    limits, range of motion).
  \item Software system integrity check (AI model version, parameter
    verification, hash validation).
  \item Network connectivity and MCP-PAI server communication test.
  \item Emergency stop functionality test.
  \item Sensor calibration verification (cameras, force sensors, proximity
    sensors).
  \end{enumerate}

\item Systems that fail any daily verification check shall be placed in
  maintenance status and shall not be used for patient procedures until
  the issue is resolved and the check is passed.

\item Daily check results shall be recorded in the system's audit trail with
  timestamp and pass or fail status for each verification element.
\end{enumerate}

\subsection{System Utilization Management}

\begin{enumerate}[label=(\alph*)]
\item The site's scheduling system shall manage Physical AI system utilization
  to balance patient demand against system availability, following the
  utilization patterns validated in the simulation:
  \begin{enumerate}[label=(\arabic*)]
  \item RT positioning robots: 38 percent average utilization.
  \item Imaging assistants: 38 percent average utilization.
  \item RT motion-tracking robots: 37 percent average utilization.
  \item Surgical robots: scheduled with appropriate intervals between
    procedures for decontamination and preparation.
  \item Companion robots and humanoids: flexible scheduling to accommodate
    patient emotional support needs~\cite{newtrial2026}.
  \end{enumerate}

\item When utilization approaches 80 percent for any system type, the
  scheduling system shall extend appointment intervals to prevent queue
  buildup and maintain the target 8-minute average wait time.

\item Maintenance windows shall be scheduled during periods of historically
  low demand (typically 22:00 to 06:00), using the simulation data as the
  baseline demand model~\cite{newtrial2026}.
\end{enumerate}

\section{Quality Assurance}

\subsection{Continuous Quality Monitoring}

\begin{enumerate}[label=(\alph*)]
\item The site shall maintain a quality management system compliant with the
  adapted ICH E6(R3) requirements for Physical AI
  trials~\cite{iche6r3adapt}.

\item Key quality indicators shall be tracked continuously:
  \begin{enumerate}[label=(\arabic*)]
  \item Patient throughput versus the 168-patient 24-hour baseline.
  \item Adverse event rate versus the 4.2 percent baseline (7 events across
    168 patients)~\cite{newtrial2026}.
  \item Same-hour adverse event resolution rate versus the 100 percent
    baseline.
  \item Fleet uptime versus the 99.7 percent baseline.
  \item Patient satisfaction versus the 4.7 out of 5.0 baseline.
  \item Cumulative PSL score versus the 63.4 to 64.4 baseline range.
  \end{enumerate}

\item Deviations from baseline quality indicators shall trigger root cause
  analysis and corrective action within 7 days.
\end{enumerate}

\subsection{Regulatory Compliance Monitoring}

\begin{enumerate}[label=(\alph*)]
\item Compliance with the three adapted regulatory frameworks shall be
  monitored continuously:
  \begin{enumerate}[label=(\arabic*)]
  \item ICH E6(R3): GCP compliance for Physical AI systems, data governance,
    audit trail integrity~\cite{iche6r3adapt}.
  \item 21 CFR Part 50: Informed consent completeness, patient rights
    adherence, pediatric protections~\cite{cfr50adapt}.
  \item 21 CFR Part 312: IND compliance, safety reporting timelines, annual
    report preparation~\cite{cfr312adapt}.
  \end{enumerate}

\item Monthly compliance reports shall be submitted to CDPH as required by
  the state regulations.

\item The Physical AI Safety Review Committee shall meet at intervals not
  exceeding 90 days to review all safety data, quality metrics, and
  compliance status~\cite{cfr312adapt}.
\end{enumerate}

\section{Adverse Event Management}

\begin{enumerate}[label=(\alph*)]
\item Adverse event detection, documentation, and initial response shall be
  automated by Physical AI systems, with site safety officer notification
  within 1 minute of detection.

\item The seven adverse event categories defined in the adapted 21 CFR Part 312
  shall be used for classification: serious Physical AI adverse events,
  Physical AI-drug interactions, cybersecurity incidents, AI model degradation,
  sim-to-real divergence, digital twin discrepancy, and safety committee
  findings~\cite{cfr312adapt}.

\item The site shall target same-hour resolution for all Grade 1 and Grade 2
  events, matching the simulation performance where all seven adverse events
  (nausea Grade 1, hypotension Grade 1, bleeding Grade 1, pain Grade 2,
  cough Grade 1, desaturation Grade 1, anxiety Grade 1) were resolved within
  the hour of occurrence~\cite{newtrial2026}.

\item Grade 3 or higher events shall trigger immediate escalation to the
  on-call emergency physician and the Physical AI Safety Review Committee.

\item All adverse events shall be reported to CDPH and FDA within the timelines
  specified in the adapted 21 CFR Part 312 \S~312.32 (7 days for fatal or
  life-threatening, 15 days for serious)~\cite{cfr312adapt}.
\end{enumerate}

